\documentclass{report} 
\title{Bloch}
\date{Started 15th June 2026}
\author{Malcolm}
\usepackage{amsmath} %import math
\usepackage{mathtools} %more math
\usepackage{amssymb} %for QED symbol
\usepackage{amsthm} %
\usepackage{mathrsfs} %
\usepackage{bm}%bold math
\usepackage{graphicx} %import imaging

\graphicspath{{./images/}} %set imaging path

\newtheorem{theo}{Theorem}[chapter]
\newtheorem{theorem}[theo]{Theorem}
\newtheorem{lemma}[theo]{Lemma}
\newtheorem{corollary}[theo]{Corollary}
\newtheorem{definition}[theo]{Definition}


\begin{document}
\maketitle

\tableofcontents

\chapter{Construction of the Real Numbers}
\begin{definition}
Let $S$ be a set. A \textbf{binary operation} on $S$ is a function $S\times S\to S$. A
\textbf{unary operation} on $S$ is a function $S\to S$.
\end{definition}
\noindent Let $S$ be a set and let $*:S\times S\to S$ be a binary operation. If $x,y\in S$, the correct way to write the result of the operation to the pair $\langle x,y\rangle$ would be
$*(\langle x,y\rangle)$. However since this is a bit cumbersome one may write $x*y$ instead.
Similarly if $\neg:S\to S$ is a unary operation and $x\in S$ one may choose to write $\neg x$ over $\neg(x)$.

\section{Axioms for the Natural numbers}
\subsection{Peano Postulates}
\textit{Axiom: Peano Postulates}\\
There exists a set $\mathbb N$ with an element $1\in\mathbb N$ and a function $s:\mathbb N\to\mathbb N$ that satisfy the following three properties.
\begin{itemize}
\item There is no $n\in\mathbb N$ such that $s(n)=1$.
\item The function $s$ is injective.
\item Let $G\subseteq\mathbb N$ be a set. Suppose that $1\in G$, and that if $g\in G$ then $s(g)\in G$. Then $G=\mathbb N$.
\end{itemize}
It will be shown that the set $\mathbb N$ is unique. We can therefore give it a name.
\begin{definition}
The set of \textbf{natural numbers} denoted $\mathbb N$, is the set the existence of which is given in the Peano Postulates.
\end{definition}
\newpage

\subsection{Consequences of Peano Postulates}
\begin{lemma} 
Let $a\in\mathbb N$. Suppose that $a\neq 1$. Then there is a unique $b\in\mathbb N$ such that $a=s(b)$. (where existence of $s$ and $\mathbb N$ are outlined by the peano axioms)
\end{lemma}
\begin{proof}
First we prove existence. Let 
\begin{equation*}
G=\{1\}\cup\{c\in\mathbb N\mid \exists b\in\mathbb N(s(b)=c)\}
\end{equation*}
(which exists by the subset, pairing, and union axioms) We will prove that $G=\mathbb N$. To do this we need to show $1\in G$, $G\subseteq\mathbb N$ and $\forall g\in G(s(g)\in G)$. 
The first two are 
straightforward. For the third suppose $g\in G$; we have two cases. First suppose $g\in\{1\}$; then $g=1\in\mathbb N$. Since $s:\mathbb N\to\mathbb N$, we then have $s(1)\in\mathbb N$.
So by definition $s(1)\in\{c\in\mathbb N\mid \exists b\in\mathbb N(s(b)=c)\}$ and $s(g)=s(1)\in G$.
Now suppose $g\in\{c\in\mathbb N\mid \exists b\in\mathbb N(s(b)=c)\}$; then by definition $g\in\mathbb N$ and $s(g)\in\mathbb N$ exists. 
By definition $s(g)\in\{c\in\mathbb N\mid \exists b\in\mathbb N(s(b)=c)\}$ so $s(g)\in G$. By the third peano postulate we then can assert that $G=\mathbb N$.
So $a\in\mathbb N=G$. Since $a\neq 1$ we must then have $a\in\{c\in\mathbb N\mid \exists b\in\mathbb N(s(b)=c)\}$ and by definition, $\exists b\in\mathbb N(s(b)=a)$.\\
\indent Now for uniqueness. Let $m$ and $n$ be arbitrary elements of $\mathbb N$ such that $a=s(m)$ and $a=s(n)$. Then $\langle m,a\rangle\in s$ and $\langle n,a\rangle\in s$. Since $s$ is injective
we must have $n=m$.
\end{proof}

\begin{theorem}

\end{theorem}









\end{document}
